\documentclass{exam}
\usepackage{verbatim}
\usepackage{xurl}
\usepackage{hyperref}
\usepackage{color}
\usepackage{tabularx}
\usepackage[utf8]{inputenc}
\usepackage{longtable}
\author{Dr. Rebin Abdulkader Muhammad}
\title{DATA 110 -- DATA VISUALIZATION AND COMMUNICATION \\(CRN 21802)\\Syllabus -- Fall 2026}

\date{}
%\pagenumbering{arabic} % sets page numbering to arabic numerals
\begin{document}

\maketitle
\begin{center}
\section{General Information}
\end{center}
\pagenumbering{arabic} % sets page numbering to arabic numerals

\subsection*{Instructor Information}
\textbf{Name:} Dr. Rebin Abdulkader Muhammad \\
\textbf{E-mail:} \href{mailto:Rebin.Muhammad@montgomerycollege.edu}{Rebin.Muhammad@montgomerycollege.edu} \\
\textbf{Office Phone Number:} 240-567-1931

\subsection*{Office Hours}
\textbf{Every day, 12:00--1:00 PM via Microsoft Teams.} Students may also contact me by e-mail or a Teams message to arrange a meeting at another mutually convenient time.

\subsection*{Class Meeting Times and Room}

\begin{center}
\begin{tabular}{| c | c | c |}
\hline
\textbf{Day} & \textbf{Time} & \textbf{Location} \\
\hline
Wednesday & 6:00 PM--8:55 PM & Microsoft Teams \\
\hline
\end{tabular}
\end{center}

\noindent \textbf{Course Information:} DATA 110, CRN 21802, 3 credits, Fall 2026. The official section dates are August 31, 2026 through December 20, 2026. The first scheduled class meeting is Wednesday, September 2, 2026, and the last scheduled class meeting is Wednesday, December 16, 2026.

\noindent \textbf{Course Modality:} This is a \textbf{synchronous online Class/Lecture} course. Class meetings are held live via \textbf{Microsoft Teams} every Wednesday from 6:00 PM--8:55 PM. Students should use the course Teams link below to attend class.

\noindent \textbf{Online Class Access:} Blackboard will be used for course materials, announcements, assignments, and other course information. The Microsoft Teams course space is available through the \href{https://teams.microsoft.com/l/team/19\%3AVSRlPh8qbhrTkXDovKRTuRAOkB_01cWJJ1WZeneAYbI1\%40thread.tacv2/conversations?groupId=7627d91b-7652-4611-947b-8faf8c4080ee\&tenantId=4d401e0b-6552-4e08-bd38-a8d15ca0ddd9}{\textcolor{blue}{General | Data110-Fall2026}} channel.

\section*{Course Information}

\subsection*{Course Description}
Emphasis on data visualization and communication skills for professional situations, including effective quantitative summary and public speaking. Preparing and producing data visualizations, presentations, and technical documents for specific audiences and analyses for general audiences.

\noindent \textbf{Credit Hours:} 3 semester hours. Three hours each week.

\subsection*{Prerequisites}
A grade of C or better in MATH 117/MATH 117A, MATH 217, BSAD 210, or consent of department.

\subsection*{Learning Outcomes}
Upon course completion, students will be able to:
\begin{enumerate}
    \item Demonstrate proficiency with methods, tools, and techniques for data visualization.
    \item Construct oral and written arguments utilizing quantitative data.
    \item Express findings from scientific data analysis proficiently to a target audience.
    \item Apply techniques to develop and relate compelling stories with data.
    \item Compose and modify analytical summaries.
    \item Describe impression management strategies and situational communication preferences.
\end{enumerate}

\subsection*{Course Materials}
\begin{itemize}
    \item \textbf{Main Textbook:} Claus O. Wilke, \textit{Fundamentals of Data Visualization}, 1st edition (2019). The complete online version is available \textbf{free of charge (\$0 cost to students)} at \href{https://clauswilke.com/dataviz/}{\textcolor{blue}{Fundamentals of Data Visualization}}.
    \item \textbf{Tools:} Python, Jupyter Notebook and/or Google Colab, Tableau, and GitHub, as assigned.
\end{itemize}

\subsection*{Course Format}
This class is highly interactive and project-based. You will work in teams to analyze datasets, create visualizations, and present findings. Weekly assignments will focus on reinforcing concepts and applying skills learned in class.

\subsection*{Course Requirements}
\begin{itemize}
    \item Regular assignments involving data analysis, visualization, and quantitative summaries.
    \item Oral presentations and technical documents designed for specific audiences.
    \item Major projects using real datasets to develop visualizations, analytical summaries, and data-driven stories.
    \item Regular classwork using programming, visualization, and communication tools.
\end{itemize}

\subsection*{Technology, Calculator, and Other Materials}
Students need regular access to a computer with a reliable internet connection and must be able to access Blackboard. Course work will use Python-based tools, Jupyter Notebook and/or Google Colab, Tableau, GitHub, and other data-visualization tools introduced during the semester.

\noindent \textbf{Webcam and Microphone:} Because this course meets synchronously through Microsoft Teams, students need access to a microphone for class participation and presentations. A webcam is also expected for interactive class activities and presentations unless the instructor indicates otherwise or an approved accommodation applies.

\noindent \textbf{Computational Tools:} Python and other course-approved software will be used for calculations, data analysis, and visualization. Any restrictions on the use of software, websites, AI tools, or other resources during a quiz or assessment will be clearly communicated in advance.

\noindent \textbf{Other Supplies:} No additional required supplies are currently anticipated beyond access to the textbook and the technology/software listed above. If an additional material becomes necessary, students will be notified in advance.

\subsection*{Assignment Submission and Collection}
Unless otherwise announced, assignments, quizzes, projects, and other required work will be submitted electronically through Blackboard or through the course platform specified in the assignment instructions. Due dates, submission locations, file-format requirements, and any special collection procedures will be stated with each assignment. Students are responsible for confirming that their work has been successfully submitted.

\subsection*{Late Work and Submission Policy}
Assignments and projects must be submitted by the posted deadline. Late submissions will not be accepted. When an activity allows multiple attempts or revisions, students are encouraged to begin early, use feedback, and resubmit before the deadline. If an unexpected or documented circumstance affects your ability to meet a deadline, please contact the instructor as soon as possible.

\subsection*{Make-Up Assessments}
Make-up quizzes, presentations, or other scheduled assessments are not automatically granted. If an unexpected, documented, or College-approved circumstance prevents you from completing an assessment as scheduled, contact the instructor as soon as possible. When a make-up is approved, the instructor will determine the appropriate format and deadline. Make-up work may differ from the original assessment while measuring the same course outcomes.

\subsection*{Python Projects}
\begin{itemize}
    \item Analyze the \textbf{Titanic Dataset} to explore survival factors using histograms and KDE plots.
    \item Visualize the \textbf{Happiness Index Dataset} to uncover global trends.
    \item Create interactive dashboards for custom datasets using \textbf{Plotly}.
    \item Perform exploratory data analysis on the \textbf{Penguins Dataset} using Python.
\end{itemize}

\subsection*{Grading}
\begin{table}[h]
    \centering
    \renewcommand{\arraystretch}{1.2}
    \begin{tabular}{|l|r|}
    \hline
    \textbf{Type of Activity} & \textbf{Percent} \\
    \hline
    Homework & 20\% \\
    \hline
    Quizzes & 10\% \\
    \hline
    Class Presentations & 10\% \\
    \hline
    Class Discussion / Participation & 10\% \\
    \hline
    Project 1 & 15\% \\
    \hline
    Final Project & 25\% \\
    \hline
    Museum Visit & 5\% \\
    \hline
    Infographic & 5\% \\
    \hline
    \textbf{Total} & \textbf{100\%} \\
    \hline
    \end{tabular}
    \caption{Grading Breakdown}
\end{table}

\section*{Course Policies}

\subsection*{Attendance}
Regular attendance and active participation are essential for success in this course. Missing class can affect your ability to complete team projects and understand course material.

\subsection*{Academic Integrity}
Academic integrity is critical for students and instructors alike. Students are expected to adhere to the \href{https://www.montgomerycollege.edu/_documents/policies-and-procedures/42001-student-code-of-conduct.pdf}{\textcolor{blue}{Montgomery College Student Code of Conduct}}. All submitted work must represent your own work or your group's work when collaboration is explicitly assigned.

\noindent \textbf{Generative Artificial Intelligence (AI):} Unless the instructor explicitly permits AI for a particular activity, submitted work must reflect the student's own thinking, analysis, and authentic voice and must be written in the student's own words. AI tools may not be used to generate, rewrite, or substantially complete work that is presented as the student's own. When AI use is explicitly permitted, students must follow the instructor's directions for acceptable use, disclosure, and citation. Students remain responsible for the accuracy, quality, and originality of everything they submit. Unauthorized or undisclosed AI-generated work may be treated as a violation of academic integrity.

\subsection*{Syllabus Changes}
While I try to stick to the plans laid out here, this syllabus is subject to change. Any updates will be communicated via Blackboard.

\section*{Important Dates}

\subsection*{DATA 110.21802 Course Dates}
\begin{itemize}
    \item \textbf{Official section start date:} August 31, 2026
    \item \textbf{First scheduled class meeting:} Wednesday, September 2, 2026
    \item \textbf{Refund Date:} September 8, 2026
    \item \textbf{No Grade / Change to Audit Deadline:} September 22, 2026
    \item \textbf{Last Date to Withdraw:} November 17, 2026
    \item \textbf{Fall Break:} November 25--29, 2026
    \item \textbf{Final Exam Week:} December 14--20, 2026
    \item \textbf{Last scheduled class meeting:} Wednesday, December 16, 2026
    \item \textbf{Official section end date:} December 20, 2026
\end{itemize}

It is the student's responsibility to officially drop or withdraw from the course. Non-attendance does not constitute official withdrawal.

For the full College calendar, see the \href{https://www.montgomerycollege.edu/_documents/academics/academic-calendar-2026-2027.pdf}{\textcolor{blue}{Academic Year Calendar 2026--2027}}.

\section*{Tips for Success}
\begin{itemize}
    \item Be an active participant during class time. Ask questions - during class, office hours, or via email.
    \item Give yourself plenty of time to prepare for speeches and projects.
    \item Study with classmates and make use of learning centers.
    \item Engage in discussions and respect diverse viewpoints.
\end{itemize}

\section*{Expectations on Attendance and Participation}

Learning shouldn't be done in isolation; it is a social activity. Attend all class meetings and please be on time. This way you will avoid missing important information and also show respect for your peers.

Being present, active, and engaged in class will provide you with the deepest learning experience. To get the most out of this class, you will need to actively engage with the materials, activities, discussions, and course resources.

Strong participation means having the assigned homework done (or attempted) on time and having questions if there are things you don't understand. It means being actively involved in discussions, asking questions, and demonstrating that you read and have thought about the material. When not in class, study with classmates and in the learning center.

Participation means demonstrating respect for others' ideas through acknowledging their views and asking for clarification when you aren't sure. Participation also means focusing on what is going on (being present), stepping up when you have a contribution and stepping back when it is time for others to talk. Whether you are outgoing or shy, everyone has something to contribute and together we will create a space for you to share.

In cases involving excessive absences from class, I may drop you from the class. Excessive absence is defined as one more absence than the number of classes per week during a fall or spring semester – in our case this means \textbf{2 absences}. Excessive absences can significantly decrease the likelihood that you will be successful in this class and it is not in your best interests to remain enrolled in a class where this is the case. I am aware that life happens. If you find yourself dealing with...

\section*{Inclusivity and Commitment to Student Success}

\begin{itemize}
    \item Your well-being and success in this course matters to me.
    \item I will actively listen to you in this course, because your voice matters to me.
    \item I expect you to actively listen to me and each of your peers in our classroom.
    \item Your unique experience and background will be accepted in our classroom.
    \item I expect you to participate in this course without fear of making mistakes!
    \item I am committed to ensuring our classroom is an inclusive learning environment.
    \item It is my intent that every student is well served by this course, that the learning needs of every student are addressed both in and out of class, and that the diversity that students bring to this class be viewed as a resource, strength, and benefit within our classroom as well as in the broader community.
    \item It is my intent to present materials and activities in a safe space that are respectful of diversity: religion, gender, sexuality, psychological experience, ability/disability, age, socioeconomic status, ethnicity, race, language, immigration status, ideological perspectives, and culture. Your suggestions are encouraged and appreciated. Please share with me ways to improve the effectiveness of this course for you personally or for other students or student groups. In addition, if any of the assignment due dates conflict with your religious observances, please let me know so that we can make arrangements.
\end{itemize}

\section*{Tools and Technologies}

For our Data Visualization class, we'll be utilizing:

\begin{itemize}
    \item \textbf{Python}: Essential for data manipulation and creating complex visualizations.
    \item \textbf{Tableau}: For crafting interactive dashboards that bring data to life.
    \item \textbf{GitHub}: To manage and share our visualization codes and projects.
\end{itemize}

These tools are chosen to give you a broad exposure to the field of data visualization.

\section*{Data Visualization Class -- Weekly Breakdown}

This tentative schedule is subject to change as needed. Updates will be communicated through Blackboard.

\begin{longtable}{|l|l|p{10cm}|}
\hline
\textbf{Date} & \textbf{Week} & \textbf{Topic} \\ \hline
September 2  & 1  & Introduction to Tools: Google Forms, Sheets, GitHub, and Markdown for documentation. \\ \hline
September 9  & 2  & Google Sheets \& Python Basics: Google Colab, Python variables, loops, and data types. \\ \hline
September 16 & 3  & Scatter Plots with Matplotlib: line/scatter plots and customization. \\ \hline
September 23 & 4  & Bar Graphs with Matplotlib: simple, grouped, and stacked bar charts. \\ \hline
September 30 & 5  & Histograms and KDE: visualizing distributions. \\ \hline
October 7    & 6  & Ethics \& Advanced Distributions: ethical data practices and complex visualizations. \\ \hline
October 14   & 7  & Penguins Dataset Analysis: exploratory data visualization using Python. \\ \hline
October 21   & 8  & Midterm Project Development: planning, analysis, and drafting. \\ \hline
October 28   & 9  & Midterm Project Presentations: present and refine projects. \\ \hline
November 4   & 10 & Visualizing Maps with Plotly: interactive geographic data visualization. \\ \hline
November 11  & 11 & Treemaps \& Advanced Maps: advanced Plotly visualization techniques. \\ \hline
November 18  & 12 & Storytelling with Data: designing audience-focused narratives with data. \\ \hline
November 25  & -- & \textbf{Fall Break -- College Closed; No Class.} \\ \hline
December 2   & 13 & Final Project Development: develop and refine final projects. \\ \hline
December 9   & 14 & Final Project Development, peer feedback, and presentation preparation. \\ \hline
December 16  & 15 & Final Project Presentations and course wrap-up. \\ \hline
\end{longtable}

\section*{Student Resources}

\subsection*{Assessment and Testing Centers}
Visit the \href{https://www.montgomerycollege.edu/admissions-registration/assessment-testing-centers/}{\textcolor{blue}{Assessment and Testing Centers}} website for current testing information and resources.

\subsection*{Disability Support Services}
Students who may need an accommodation due to a disability should contact Disability Support Services (DSS). Approved accommodations will be provided in accordance with the DSS accommodation letter.
\url{https://www.montgomerycollege.edu/counseling-and-advising/disability-support-services/index.html}

\subsection*{Counseling \& Advising}
\url{https://www.montgomerycollege.edu/counseling-and-advising/}

\subsection*{Learning Centers and Academic Support Centers}
\url{https://www.montgomerycollege.edu/academics/support/learning-centers/index.html}

\subsection*{Student Forms}
\url{https://www.montgomerycollege.edu/admissions-registration/student-resources/student-forms.html}

\subsection*{Safety, Security, \& Emergency Operations}
\url{https://www.montgomerycollege.edu/life-at-mc/public-safety/index.html}

\subsection*{Pregnant and Parenting Students}
Montgomery College provides Title IX protections and academic support for students who are pregnant or experiencing pregnancy-related conditions. Medically necessary absences will be excused and missed work may be made up as required by College policy. See \href{https://www.montgomerycollege.edu/policies-and-procedures/title-ix/pregnant-parenting-students.html}{\textcolor{blue}{Pregnant and Parenting Students}} for details.

\section*{Main Repository}
The course repository for \textbf{DATA 110 (CRN 21802)} will be added here when it is available.

\end{document}
